% Post-CP21 higher-degree intermediate; no irreducibility assertion.
\subsection{Actual natural blocks in every odd degree on a fixed nonordinary surface}
\status{Proved and independently checked.}

\begin{lemma}[Sharp pure norms in every odd principal degree]
\label{lem:nonordinary-ppav-odd-sharp-norm-candidate}
Let $(A,\Theta)$ be a principally polarized abelian surface over an
algebraically closed field of characteristic two, with two-rank $f<2$.
For every odd positive integer $d$, there is an odd polarized isogeny
$\phi:(B,\Theta_B)\to(A,\Theta)$ of degree $d^2$ and multiplier $d$.
After one constant target translation, the norms of a nonempty open
of smooth maximal-rank curves $C\in|3\Theta_B|$ form integral nodal
curves $D\in|3d\Theta|$ with normalization $C$ and
\[
             \delta(D)=9(d^2-1),\qquad
             f(D)-f=9d^2-1.
\]
Thus this entire open norm family lies in the maximal-relative-unit-rank
locus on the target. Its graph representation has image a subgroup
of a product of symmetric groups $S_{18}$.
\end{lemma}
\begin{proof}
Choose a maximal isotropic subgroup of $A[d]$, take its principal
polarized quotient, and use the polarized dual isogeny. The standard
finite theta-group descent gives the stated $B,\phi$ and
$\phi^*\lambda_\Theta=d\lambda_{\Theta_B}$; this is the same
odd principal-isogeny construction used in
Theorem~\ref{thm:DE-A3-ppav-ranks}. Its degree is $d^2$ and it is
\etale. In particular $B$ has the same two-rank $f$ as $A$.
Put $K=\ker\phi$, a constant group of odd order, and $L_B=3\Theta_B$.
The smooth maximal-rank open in $|L_B|$ is nonempty and, by
Theorem~\ref{thm:every-nonordinary-ppav-degree-one-full-candidate},
has relative image $\GL_8(\mathbf F_2)$.

We detail the nodality input, which uses no ordinarity of the ambient
surface. A general member is invariant under no nonzero $g\in K$:
if $g$ does not preserve $L_B$, invariance is impossible; otherwise
it acts nontrivially on the complete very ample projective system,
so its invariant members lie in a proper union of eigenspaces.
On a nonproduct source, for each fixed $g\ne0$ the two first jets
at $P,P+g$ have rank six except at finitely many $P$, by
Corollary~\ref{cor:all-theta-fixed-point-jets-candidate}.
Indeed its exceptional set is contained in
\[
                     3P+g\in\Theta_B,\qquad
                     3P+2g\in\Theta_B.
\]
The two theta translates are distinct and intersect finitely, and
$[3]$ is finite. On the rank-six open, the two value equations and
proportionality of the nonzero gradients impose three conditions;
$P$ varies in dimension two, so tangency is absent for a general
member. The finite exceptional set already imposes two independent
value conditions and is avoided. Three distinct points impose three
independent values on $3\Theta_B$, by the theta-product separation
argument in Lemma~\ref{lem:ordinary-ppav-isogeny-block}. Applying this
to $P,P+g,P+h$ excludes triple identifications, since their point
parameter has dimension two.

For completeness the product case uses only the cubic-product jet
calculation, not ordinary factors. If the two translated points differ
in both coordinates, Lemma~\ref{lem:product-small-degree-jets} says
that failure of rank six requires the cross-tangency equations
$3P_1+2g_1=0$, $3P_2+g_2=0$, or the equations with $g$ and $2g$
interchanged. This gives only finitely many points, since $[3]$
is finite, and two independent values avoid them. If one coordinate
agrees, the joint jets have rank exactly five everywhere: two
distinct points of a smooth plane cubic have distinct tangent lines,
as otherwise their line would have intersection multiplicity at least
four with the cubic. In the tensor first-jet calculation the two
three-dimensional jet spaces therefore meet in exactly one dimension.
After imposing the two values, the remaining three-dimensional space
of pairs of gradients cannot consist entirely of proportional pairs:
a linear subspace of rank-at-most-one $2$-by-$2$ matrices has dimension
at most two. Thus tangency again adds one condition, with no additional
axis exception. These statements hold for both ordinary and
supersingular plane cubics. Three value vectors of the Segre cubic
product are independent unless one coordinate is shared and the
three vectors in the other plane cubic are collinear. For a fixed
triple of odd translations this collinearity requires
$3P_1+g_1+h_1=0$ (or its second-factor version), so that exceptional
point locus has dimension one; two independent values still exclude
it generically. This proves the same nodality statement on a product.

There are finitely many kernel elements and pairs, so all these
nonempty conditions hold simultaneously. Thus $C\to D=\phi(C)$ is
birational, $C$ is the normalization, and $D$ is integral nodal.
For each unordered type $\{g,-g\}$, the intersections
$C\cap(C+g)$ give 18 distinct nodes, since $L_B^2=18$ and triple
identifications are excluded. Hence the number of nodes is
$18(|K|-1)/2=9(d^2-1)$. Their permutation image is a subgroup of
$\prod_{\{g,-g\}}S_{18}$.

Choose symmetric representatives of the two principal polarizations.
The norm of $L_B$ has N\'eron--Severi class $3d\Theta$ and is symmetric;
its difference from that precise target bundle is geometric
$2$-torsion in $\Pic^0(A)$. Since $3d$ is odd, translation by a
constant geometric $2$-torsion point corrects this difference,
as in Lemma~\ref{lem:ordinary-ppav-isogeny-block}.
This identifies the norm family with a family in the exact system
$|3d\Theta|$.

The normalization has genus ten and two-rank $8+f$.
Graph cohomology of an integral nodal curve adds one unit direction
per node, so
\[
             f(D)=8+f+9(d^2-1),\qquad
             f(D)-f=9d^2-1.
\]
This equals the full relative coherent dimension in the odd target
degree, and therefore proves maximal rank there. Notice that no
additional ample component has been attached to the norm.
\end{proof}

\begin{theorem}[An actual eight-dimensional block in every odd degree]
\label{thm:nonordinary-ppav-odd-natural-block-candidate}
Let $(A,\Theta)$ be any nonordinary principally polarized abelian
surface over a perfect field of characteristic two. For every odd
$d\ge1$, the degree-$d$ geometric relative monodromy contains a
subgroup conjugate to
\[
                      \operatorname{diag}(\GL_8(\mathbf F_2),I)
                  \subset\GL_{9d^2-1}(\mathbf F_2).
\]
In particular this is an actual natural block fixing a complementary
summand pointwise, rather than merely a general-linear quotient.
\end{theorem}
\begin{proof}
All constructions may be made after passing to an algebraic closure,
which is sufficient for the geometric subgroup assertion. For $d=1$
the preceding degree-one theorem gives the claim. For odd $d>1$,
use the norm family of the lemma and shrink its source open as needed.
Such shrinking does not change the full degree-one geometric image.
The normalization sequence, together with the odd-isogeny isomorphism
\[
             H^1_{\et}(A,\mathbf F_2)
                      \xrightarrow{\sim}H^1_{\et}(B,\mathbf F_2),
\]
gives the exact sequence of relative local systems
\[
 0\longrightarrow T\longrightarrow V_D\longrightarrow U_C
       \longrightarrow0,\qquad
 \dim T=9(d^2-1),\quad \dim U_C=8.
\]
Here $T$ is the graph permutation representation and the image on
$U_C$ is $\GL_8(\mathbf F_2)$. The image on $T$ has order prime to31,
since it embeds in a product of $S_{18}$, whereas31 divides
$|\GL_8(\mathbf F_2)|$. The latter group is simple.
Goursat's lemma therefore implies that the kernel of the action on
$T$ still surjects onto $\GL(U_C)$: any common nontrivial quotient
would have to be that simple group, which cannot be a quotient of
a group of order prime to31.

Choose a vector-space splitting of the displayed sequence and restrict
to this kernel. Its image lies in
$\Hom(U_C,T)\rtimes\GL(U_C)$. Its translation kernel is a
$\GL(U_C)$-submodule of $U_C^*\otimes T$, hence equals
$U_C^*\otimes T'$ for a subspace $T'\subset T$. This assertion follows
from the full matrix algebra generated by the natural general-linear
representation; it does not require a semisimplicity assertion for
arbitrary modules. Modulo $T'$ the extension is the graph of a
one-cocycle. The elementary vanishing
\[
                 H^1(\GL_8(\mathbf F_2),U_C^*)=0
\]
of Lemma~\ref{lem:char2-group-inputs} makes that cocycle a coboundary.
After conjugating the splitting, every residual upper-right column
lies in the known translation kernel, so it can be removed by an
actual kernel element. Thus the image contains the claimed
block-diagonal subgroup.

Finally the nodal norm family is inside the target maximal-rank
open, by the preceding lemma. The relative Lang system extends
\etale over this open, including its nodal curves. Its monodromy
on the dense smooth-member open is the same, because the ambient
open in the projective parameter space is normal and a dense-open
inclusion induces a surjection on geometric fundamental groups.
Consequently the subgroup just obtained belongs to the geometric
monodromy requested in the theorem.
\end{proof}

\begin{remark}[The remaining odd-degree irreducibility gate]
\label{rem:nonordinary-odd-block-scope-candidate}
The actual block above gives transvections but does not prove that
the whole degree-$d$ representation is irreducible or that its image
is the full general linear group. In particular the product common-grid
argument has not been imported for a simple surface. The remaining
graph-kernel and global invariant-subspace questions are separate.
\end{remark}
